\section{Conclusion}
\label{sec:conclusion}

We presented \textit{\ours{}}, a framework that reinforces visual latent planning for vision-language-action reasoning tasks. By combining action-aligned reinforcement learning with reasoning-enhanced action adaptation, \ours{} enables embodied agents to think before acting and execute robust actions in dynamic environments. Through extensive experiments across embodied reasoning and robot manipulation benchmarks, we demonstrated strong long-horizon planning, few-shot adaptation, and emergent behaviors such as failure detection and self-correction, providing a scalable path toward more deliberative and adaptable embodied AI systems.

\paragraph{Limitations}
Since ThinkAct builds on pretrained multimodal LLMs, it inevitably inherits their limitations, particularly hallucinations in visual or spatial reasoning. This can lead to generated plans that reference incorrect object attributes or spatial relationships, affecting downstream execution. While our latent planning and action grounding mitigate this to some extent, future work on grounding-aware training or hallucination suppression in MLLMs may further improve robustness and reliability in real-world deployment.

\paragraph{Broader Impacts}
Our work aims to enhance the reasoning capabilities of embodied agents, which could support real-world applications such as assistive robotics, home automation, and industrial systems. In particular, models like ThinkAct may help robots better interpret vague instructions and execute multi-step plans in dynamic environments. However, increased autonomy and reasoning ability in embodied systems also raise potential concerns. Misinterpretation of ambiguous commands, reliance on hallucinated visual reasoning, or overconfidence in CoT outputs could result in unintended behaviors, especially in safety-critical settings. Hence, future research on safeguards or alignment with human intent could further help mitigate these risks.